\documentclass[aspectratio=169, 11pt]{beamer}

% --- Packages & Setup ---
\usepackage[utf8]{inputenc}
\usepackage[T1]{fontenc}
\usepackage{helvet}
\renewcommand{\familydefault}{\sfdefault}
\usepackage{xcolor}
\usepackage{booktabs}
\usepackage{graphicx}
\usepackage{tikz}
\usetikzlibrary{calc, arrows.meta}

% small picture inside an architecture box (variant B)
\newcommand{\archpic}[1]{\includegraphics[width=1.55cm,height=0.9cm,keepaspectratio]{#1}}

\graphicspath{{../}{./}{figures/}}

% A figure is shown when the file exists; otherwise a labelled box marks where it goes.
\newcommand{\figorbox}[3]{%
  \IfFileExists{figures/#1}{\includegraphics[width=#2]{figures/#1}}{%
    \fbox{\begin{minipage}[c][3.2cm]{#2}\centering\scriptsize\textbf{[#3]}\end{minipage}}}}

% --- Branding Colors ---
\definecolor{amritaMaroon}{HTML}{A4123F}

% --- Theme Setup ---
\usetheme{Madrid}
\usecolortheme[named=amritaMaroon]{structure}
\setbeamertemplate{navigation symbols}{}
\setbeamertemplate{footline}[frame number]

\begin{document}

% =========================================================================
% TITLE SLIDE (Application-Based Project)
% =========================================================================
\begin{frame}
    \begin{tikzpicture}[remember picture,overlay]
        \node[fill=amritaMaroon, text=white, font=\bfseries\footnotesize, anchor=north east, rounded corners=2pt, inner sep=5pt]
        at ($(current page.north east) + (-0.2cm,-0.2cm)$)
        {Type 2: Application Project};
    \end{tikzpicture}

    \centering
    \includegraphics[width=2.4cm]{logo.png}
    \vspace{0.05cm}

    {\large \textbf{\textcolor{amritaMaroon}{Autonomous First-Response Crime Scene Exploration\\Using a Lightweight Onboard Intelligence System}}} \\
    \vspace{0.05cm}
    {\footnotesize \textbf{23CSE498 -- Project Phase II (External Review)} \quad | \quad \textbf{Team Number: G040}} \\
    \vspace{0.1cm}

    \begin{table}
        \centering
        \scriptsize
        \renewcommand{\arraystretch}{1.0}
        \begin{tabular}{|c|c|l|}
            \hline
            \textbf{S.No} & \textbf{Register Number} & \textbf{Student Name} \\
            \hline
            1 & CB.SC.U4CSE23062 & Karthik Raju \\
            \hline
            2 & CB.SC.U4CSE23063 & Maheswar Raju \\
            \hline
            3 & CB.SC.U4CSE23112 & Vijay Krishna Burra \\
            \hline
            4 & CB.SC.U4CSE23310 & Dhinesh Chandra Boddupally \\
            \hline
            5 & CB.SC.U4CSE23421 & Kuchipudi Yokshith Sai Ram \\
            \hline
        \end{tabular}
    \end{table}
    \vspace{0.05cm}

    {\footnotesize Guide: \textbf{Mr. Rajeshwar Yadav}, Assistant Professor, Department of CSE} \\
    \vspace{0.05cm}
    {\scriptsize \today}
\end{frame}

% =========================================================================
% GUIDE APPROVAL EMAIL CONFIRMATION (Mandatory)
% =========================================================================
\begin{frame}{Guide Review \& Submission Approval}
    \framesubtitle{Confirmation of Project Review and Approval from Guide}
    \centering
    \fbox{\begin{minipage}[c][5.4cm]{0.94\textwidth}
        \centering
        \textbf{[Insert Screenshot of Guide Approval Email Here]} \\[0.3cm]
        {\footnotesize [Mandatory: Attach a clear, readable screenshot of the email received from your project guide approving this presentation and final review submission]}
    \end{minipage}}
\end{frame}

\section{Requirements \& Solution Design}

% --- SLIDE 1: PROBLEM DEFINITION & REQUIREMENT ANALYSIS ---
\begin{frame}{1. Problem Definition \& Requirement Analysis}
    \framesubtitle{Target Problem, User Needs \& Functional / Non-Functional Requirements}
    \begin{block}{Problem Statement \& Target Stakeholders}
        \footnotesize
        \begin{itemize}
            \item \textbf{Real-World Problem:} before entering a crime scene, hostage site or collapsed building, responders have no map, cannot see inside safely, and do not know what kind of incident it is.
            \item \textbf{Target Users \& Use Context:} police first-response and forensic teams, search-and-rescue units; a quadruped robot goes in first.
        \end{itemize}
    \end{block}

    \begin{block}{Functional Requirements}
        \footnotesize
        \begin{itemize}
            \item \textbf{FR-1:} explores an unknown place without an operator and returns to its start point with a map of the building.
            \item \textbf{FR-2:} tracks its own position and builds a floor map using one depth camera, a gyro and leg movements.
            \item \textbf{FR-3:} from camera video, detects whether a crime is happening, the kind of emergency, and which of 13 crimes it is.
        \end{itemize}
    \end{block}

    \begin{block}{Non-Functional Requirements \& Constraints}
        \footnotesize
        \begin{itemize}
            \item \textbf{NFR-1:} real time: 20 Hz control loop, median 29 ms per step on a CPU, no GPU.
            \item \textbf{NFR-2:} safety: zero collisions, recovery from every stall, same code in every environment.
        \end{itemize}
    \end{block}
\end{frame}

% --- SLIDE 2: END-TO-END APPLICATION ARCHITECTURE ---
% Two versions of the same diagram: A (text only) and B (pictures in the boxes). Keep one, delete the other.
% Shared styles for both versions.
\tikzset{
    archbox/.style={draw=amritaMaroon, thick, rounded corners=3pt, fill=amritaMaroon!8, align=center,
        font=\scriptsize, text width=1.65cm, minimum height=1.25cm, inner sep=4pt},
    archbottom/.style={draw=amritaMaroon, thick, rounded corners=3pt, fill=amritaMaroon!8, align=center,
        font=\scriptsize, text width=2.4cm, minimum height=1.3cm, inner sep=5pt},
    archarr/.style={-{Latex[length=2mm]}, thick, amritaMaroon},
    archsection/.style={font=\small\bfseries, text=amritaMaroon},
    archlab/.style={font=\tiny, text=black!85, fill=white, inner sep=2pt, align=center}
}

% ---------- VERSION A: TEXT ONLY ----------
\begin{frame}{2. End-to-End Application Architecture}
\framesubtitle{System Architecture \& Sensor-to-Action Workflow}
\centering
\resizebox{\textwidth}{!}{%
\begin{tikzpicture}

% section headings
\node[archsection] at (8.15,3.35) {ROBOT ONBOARD: NAVIGATION AGENTS (CPU, 20 Hz)};
\node[archsection] at (9.2,-3.25) {SCENE UNDERSTANDING \& EVIDENCE};

% top row: sensing to legs
\node[archbox] (kin)  at (1.1,1.45)   {\textbf{Kinect + IMU}\\[3pt]depth, colour\\gyro, odometry};
\node[archbox] (per)  at (3.45,1.45)  {\textbf{Perception}\\[3pt]depth frame\\to 2D scan};
\node[archbox] (slam) at (5.8,1.45)   {\textbf{SLAM}\\[3pt]ICP pose\\loop closure};
\node[archbox] (map)  at (8.15,1.45)  {\textbf{Mapping}\\[3pt]occupancy grid\\5 cm cells};
\node[archbox] (fro)  at (10.5,1.45)  {\textbf{Exploration}\\[3pt]route waypoint\\or frontier goal};
\node[archbox] (plan) at (12.85,1.45) {\textbf{Planning}\\[3pt]A* path};
\node[archbox] (loco) at (15.2,1.45)  {\textbf{Locomotion}\\[3pt]trot gait\\12 joints};

\draw[archarr] (kin.east) -- (per.west);
\draw[archarr] (per.east) -- (slam.west);
\draw[archarr] (slam.east) -- (map.west);
\draw[archarr] (map.east) -- (fro.west);
\draw[archarr] (fro.east) -- (plan.west);
\draw[archarr] (plan.east) -- (loco.west);

% closed loop: moving changes what the sensors see
\draw[archarr] (loco.north) -- ++(0,0.6) -| node[pos=0.25, above, archlab] {robot moves: new view, leg odometry, gyro} (kin.north);

% bottom row: scene understanding and evidence
\node[archbottom] (vid)   at (3.45,-1.95) {\textbf{Video model}\\[4pt]VideoMamba-Small\\+ HorNet gated conv.};
\node[archbottom] (cls)   at (7.3,-1.95)  {\textbf{Classifier}\\[4pt]crime or not, kind of\\emergency, crime type (13)};
\node[archbottom] (log)   at (11.15,-1.95) {\textbf{Logging}\\[4pt]map log, hash\\ROI clips};
\node[archbottom] (cloud) at (15.0,-1.95) {\textbf{Cloud sync}\\[4pt]upload for\\responders};

\draw[archarr] (vid.east) -- node[above=5pt, archlab] {clip scores} (cls.west);
\draw[archarr] (cls.east) -- node[above=5pt, archlab] {crime, time} (log.west);
\draw[archarr] (log.east) -- node[above=5pt, archlab] {upload} (cloud.west);

% colour video from the Kinect; map and pose to the log
\draw[archarr] (kin.south) -- ++(0,-0.99) -| node[pos=0.25, above, archlab] {colour video} (vid.north);
\draw[archarr] (map.south) -- ++(0,-0.99) -| node[pos=0.25, above, archlab] {map, pose} (log.north);
\end{tikzpicture}}

\vspace{0.15cm}
{\scriptsize Agents never call each other: they exchange the latest messages on one publish--subscribe bus, each at its own rate; a supervisor keeps the mission state and a graph of rooms. After exploration the robot returns to its start point.}
\end{frame}

% ---------- VERSION B: PICTURES IN THE BOXES ----------
\begin{frame}{2. End-to-End Application Architecture}
\framesubtitle{System Architecture \& Sensor-to-Action Workflow}
\centering
\resizebox{\textwidth}{!}{%
\begin{tikzpicture}[
    archbox/.append style={minimum height=2.15cm},
    archbottom/.append style={minimum height=2.15cm}
]

% section headings
\node[archsection] at (8.15,3.85) {ROBOT ONBOARD: NAVIGATION AGENTS (CPU, 20 Hz)};
\node[archsection] at (9.2,-3.75) {SCENE UNDERSTANDING \& EVIDENCE};

% top row: sensing to legs
\node[archbox] (kin)  at (1.1,1.6)   {\textbf{Kinect + IMU}\\[2pt]\archpic{img_rgbd.png}\\[1pt]depth, colour, gyro};
\node[archbox] (per)  at (3.45,1.6)  {\textbf{Perception}\\[2pt]\archpic{icon_scan.png}\\[1pt]2D scan};
\node[archbox] (slam) at (5.8,1.6)   {\textbf{SLAM}\\[2pt]\archpic{img_map2.png}\\[1pt]ICP pose};
\node[archbox] (map)  at (8.15,1.6)  {\textbf{Mapping}\\[2pt]\archpic{icon_grid.png}\\[1pt]occupancy grid};
\node[archbox] (fro)  at (10.5,1.6)  {\textbf{Exploration}\\[2pt]\archpic{img_map.png}\\[1pt]frontier goal};
\node[archbox] (plan) at (12.85,1.6) {\textbf{Planning}\\[2pt]\archpic{icon_path.png}\\[1pt]A* path};
\node[archbox] (loco) at (15.2,1.6)  {\textbf{Locomotion}\\[2pt]\archpic{img_robot.png}\\[1pt]trot gait};

\draw[archarr] (kin.east) -- (per.west);
\draw[archarr] (per.east) -- (slam.west);
\draw[archarr] (slam.east) -- (map.west);
\draw[archarr] (map.east) -- (fro.west);
\draw[archarr] (fro.east) -- (plan.west);
\draw[archarr] (plan.east) -- (loco.west);

% closed loop: moving changes what the sensors see
\draw[archarr] (loco.north) -- ++(0,0.6) -| node[pos=0.25, above, archlab] {robot moves: new view, leg odometry, gyro} (kin.north);

% bottom row: scene understanding and evidence
\node[archbottom] (vid)   at (3.45,-2.05) {\textbf{Video model}\\[2pt]\archpic{icon_frames.png}\\[1pt]VideoMamba + HorNet};
\node[archbottom] (cls)   at (7.3,-2.05)  {\textbf{Classifier}\\[2pt]\archpic{icon_nn.png}\\[1pt]crime or not, kind, type};
\node[archbottom] (log)   at (11.15,-2.05) {\textbf{Logging}\\[4pt]map log, hash\\ROI clips};
\node[archbottom] (cloud) at (15.0,-2.05) {\textbf{Cloud sync}\\[4pt]upload for\\responders};

\draw[archarr] (vid.east) -- node[above=5pt, archlab] {clip scores} (cls.west);
\draw[archarr] (cls.east) -- node[above=5pt, archlab] {crime, time} (log.west);
\draw[archarr] (log.east) -- node[above=5pt, archlab] {upload} (cloud.west);

% colour video from the Kinect; map and pose to the log
\draw[archarr] (kin.south) -- ++(0,-0.7) -| node[pos=0.25, above, archlab] {colour video} (vid.north);
\draw[archarr] (map.south) -- ++(0,-0.7) -| node[pos=0.25, above, archlab] {map, pose} (log.north);
\end{tikzpicture}}

\vspace{0.1cm}
{\scriptsize Agents never call each other: they exchange the latest messages on one publish--subscribe bus, each at its own rate; a supervisor keeps the mission state and a graph of rooms. After exploration the robot returns to its start point.}
\end{frame}

% --- SLIDE 2b: NAVIGATION AGENTS & MISSION FLOW ---
\begin{frame}{2. End-to-End Application Architecture (Navigation)}
    \framesubtitle{Seven Agents on One Message Bus \& One Exploration Mission}
    \begin{columns}[c]
        \column{0.62\textwidth}
        \centering
        \figorbox{fig_navigation_pipeline.png}{\linewidth}{Navigation agents: figures/fig\_navigation\_pipeline.png}
        \column{0.36\textwidth}
        \centering
        \figorbox{fig_control_flow.png}{0.72\linewidth}{Control flow: figures/fig\_control\_flow.png}
    \end{columns}
    \vspace{0.1cm}
    \begin{block}{Architecture Summary}
        \footnotesize Agents never call each other: one 20 Hz scheduler runs each at its own rate and they exchange the latest messages on the bus. A mission ends by returning to the start point.
    \end{block}
\end{frame}

% --- SLIDE 2c: SCENE-UNDERSTANDING NETWORK ---
\begin{frame}{2. End-to-End Application Architecture (Scene Understanding)}
    \framesubtitle{VideoMamba-Small with a HorNet Gated Convolution}
    \centering
    \figorbox{fig_recognition_model.png}{0.97\textwidth}{Recognition network: figures/fig\_recognition\_model.png}
    \vspace{0.1cm}
    \begin{block}{Architecture Summary}
        \footnotesize An 8-frame clip becomes spatio-temporal patches, passes 24 bidirectional Mamba layers (last 8 with a HorNet gated convolution) and is classified into 13 crimes or Background.
    \end{block}
\end{frame}

% --- SLIDE 3: ARCHITECTURAL COMPONENTS & INTERFACE DESIGN ---
\begin{frame}{3. Architectural Components \& Interface Design}
    \framesubtitle{Component Breakdown, Data Flow \& Technology Specifications}
    \begin{columns}[T]
        \column{0.50\textwidth}
        \begin{block}{Core Application Components}
            \scriptsize
            \begin{itemize}
                \item \textbf{SLAM:} point-to-line ICP against a 16-keyframe submap, gyro calibration, loop closure in a robust pose graph.
                \item \textbf{Exploration:} frontier viewpoints 0.5 m from unknown space, operator route, A* planning, narrow-gap mode.
                \item \textbf{Scene understanding:} VideoMamba-Small + HorNet gnConv; crime or not, kind of emergency, crime type.
            \end{itemize}
        \end{block}

        \begin{block}{User Interaction \& Workflow Design}
            \scriptsize
            \begin{itemize}
                \item \textbf{Primary User Journey:} load a scene, optionally draw a search route, start the mission; the robot explores and returns.
                \item \textbf{Interface Design:} live map view, per-run report and plots; one command takes a video and prints the three-level answer.
            \end{itemize}
        \end{block}

        \column{0.46\textwidth}
        \begin{block}{Data Management \& External Interfaces}
            \scriptsize
            \begin{itemize}
                \item \textbf{Data Storage:} per-run folder with metrics, trajectory, maps and the full sensor recording (replayable offline).
                \item \textbf{External Interfaces:} CoppeliaSim ZMQ API: Kinect depth and colour, gyro, 12 joints, contact check.
            \end{itemize}
        \end{block}

        \begin{block}{Platform \& Deployment Specifications}
            \scriptsize
            \begin{itemize}
                \item \textbf{Development Stack:} Python 3.12, NumPy, OpenCV, CoppeliaSim; PyTorch 2.1, mamba\_ssm for scene understanding.
                \item \textbf{Deployment Target:} navigation on a PC CPU (Ryzen 7, no GPU); scene understanding on a Tesla V100 GPU.
            \end{itemize}
        \end{block}
    \end{columns}
\end{frame}

% =========================================================================
% SECTION 2: IMPLEMENTATION COMPLETION (100%)
% =========================================================================
\section{Implementation Details}

% --- SLIDE 4: IMPLEMENTATION COMPLETION (100%) - OVERVIEW ---
\begin{frame}{4. Implementation Details -- Overview}
    \framesubtitle{End-to-End Application Module \& Integration Status}
    \begin{table}
        \centering
        \scriptsize
        \renewcommand{\arraystretch}{1.2}
        \begin{tabular}{p{3.4cm} p{7.4cm} c}
            \toprule
            \textbf{Application Module} & \textbf{Implemented Features \& Integration Scope} & \textbf{Status} \\
            \midrule
            Operator interface & Mission start, scene profile with operator route, live map, run report; video command with three-level answer & Functional \\
            SLAM \& exploration & 7 agents on the message bus: perception, SLAM, mapping, frontier, planning, locomotion, supervisor & 100\% (Integrated) \\
            Scene understanding & Trained network: crime or not, kind of emergency, crime type, strongest moment & 100\% (Functional) \\
            Data layer & Run folders, sensor recordings, offline replay, evaluation against the scene model & 100\% (Integrated) \\
            External / IO & CoppeliaSim: Kinect depth and colour, gyro, joints, contact check & 100\% (Functional, sim) \\
            \bottomrule
        \end{tabular}
    \end{table}

    \begin{block}{Prototype Environment \& Repository}
        \footnotesize
        \begin{itemize}
            \item \textbf{Execution Platform:} CoppeliaSim on Windows 11 (AMD Ryzen 7, CPU only); Linux cluster with a Tesla V100 for scene understanding.
            \item \textbf{Source Repository:} github.com/yokshith09/Autonomous-Quadruped-SLAM-Navigation-Perception (tag final-multi-env-v2); github.com/vijaykrishnaaa/VideoMamba-HorNet-Crime-Recognition
        \end{itemize}
    \end{block}
\end{frame}

% --- SLIDE 5: FUNCTIONAL MODULE IMPLEMENTATION DETAILS ---
\begin{frame}{5. Functional Module Implementation Details}
    \framesubtitle{Key Feature Implementation \& Inter-Module Integration}
    \begin{columns}[T]
        \column{0.50\textwidth}
        \begin{block}{Navigation: SLAM \& Exploration}
            \scriptsize
            \begin{itemize}
                \item \textbf{Implemented Workflow:} depth frame $\rightarrow$ 2D scan $\rightarrow$ pose (ICP + loop closure) $\rightarrow$ map $\rightarrow$ goal $\rightarrow$ A* path $\rightarrow$ trot gait; back to the start when done.
                \item \textbf{Inputs \& Outputs:} Kinect depth, gyro, leg odometry in; pose, occupancy map and joint commands out.
                \item \textbf{Validation \& Controls:} a scan match is accepted only if it agrees with odometry; watchdog recovers stalls; drop check stops at floor edges.
            \end{itemize}
        \end{block}

        \begin{block}{Scene Understanding}
            \scriptsize
            \begin{itemize}
                \item \textbf{Implemented Workflow:} a video is cut into 60 clips at four scales; the four strongest clips per crime decide the answer.
                \item \textbf{Inter-Module Integration:} reads the same Kinect colour stream; runs as its own program on a GPU.
            \end{itemize}
        \end{block}

        \column{0.46\textwidth}
        \centering
        \vspace{0.1cm}
        \figorbox{fig_scene2.png}{\linewidth}{Scene 2 map: figures/fig\_scene2.png}
        {\scriptsize Rescue field: layout (left) and the robot's map with the true (blue) and estimated (orange) path; the robot passes through the damaged building.}
    \end{columns}
\end{frame}

% =========================================================================
% SECTION 3: FUNCTIONAL TESTING & DEMONSTRATION
% =========================================================================
\section{Functional Testing \& Demonstration}

% --- SLIDE 6: FUNCTIONAL TESTING & VERIFICATION RESULTS ---
\begin{frame}{6. Functional Testing \& Verification Results}
    \framesubtitle{End-to-End Feature Verification \& Test Case Summary}
    \begin{table}
        \centering
        \scriptsize
        \renewcommand{\arraystretch}{1.15}
        \begin{tabular}{p{3.3cm} p{3.4cm} p{4.6cm} c}
            \toprule
            \textbf{Test Case / Feature} & \textbf{Test Input / Action} & \textbf{Verified Outcome} & \textbf{Status} \\
            \midrule
            TC-01: Exploration, building & Unknown 10 $\times$ 10 m building, 30 min & 94.2\% of floor, 0 collisions, ATE 14.8 cm & Pass \\
            TC-02: Route through a building & Rescue field, 70-waypoint route, 40 min & Through the building, 0 collisions, 5 stalls recovered & Pass \\
            TC-03: SLAM replay & Stored recordings replayed & Reproduces live error exactly & Pass \\
            TC-04: Boundary cases & Blocked waypoint, stall, floor edge & Skipped, recovered, stopped 0.44 m before edge & Pass \\
            TC-05: Crime or not & 290 UCF-Crime test videos & AUC 90.43\% & Measured \\
            TC-06: Crime type & 140 crime test videos, 13 classes & 43.10\% $\pm$ 1.47 & Measured \\
            \bottomrule
        \end{tabular}
    \end{table}

    \begin{block}{Application Verification Summary}
        \scriptsize
        \begin{itemize}
            \item \textbf{Functional Verification:} two environments, same code: 21.8 cm mean position error, 90.0\% floor covered, map F1 78.0\%, no collisions; 91 automated tests pass.
            \item \textbf{Responsiveness \& Operational Metrics:} 29 ms median control step on a CPU; scene understanding answers a 10--14 s video in 9--14 s on the GPU.
        \end{itemize}
    \end{block}
\end{frame}

% --- SLIDE 6b: BENCHMARKING AGAINST OTHER SYSTEMS ---
\begin{frame}{6. Functional Testing \& Verification Results (Benchmarks)}
    \framesubtitle{Same Data, Same Scoring: Where the System Stands}
    \begin{columns}[T]
        \column{0.50\textwidth}
        \begin{block}{Localisation on our recordings (ATE)}
            \scriptsize
            \begin{tabular}{lcc}
                \toprule
                \textbf{Method} & \textbf{Scene 1 / 2} & \textbf{Average} \\
                \midrule
                Odometry only & 291.4 / 15.7 cm & 153.6 cm \\
                BreezySLAM (best of 11) & 189.7 / 109.5 cm & 149.6 cm \\
                Cartographer (best of 8) & 341.5 / 85.2 cm & 213.4 cm \\
                \textbf{Ours} & \textbf{11.7 / 12.3 cm} & \textbf{12.0 cm} \\
                \bottomrule
            \end{tabular}
        \end{block}
        \begin{block}{Exploration vs nearest-frontier (same scenes)}
            \scriptsize
            Floor covered 90.0 vs 78.7\%; stuck 3 vs 12.5 per run; only ours passed through the building.
        \end{block}

        \column{0.46\textwidth}
        \begin{block}{Public TUM RGB-D Pioneer (ATE)}
            \scriptsize
            \begin{tabular}{lc}
                \toprule
                \textbf{System} & \textbf{Average} \\
                \midrule
                ORB-SLAM3 (published) & 172.0 cm \\
                OpenVSLAM (published) & 39.2 cm \\
                RTABMap (published) & 28.7 cm \\
                Wheel odometry only & 28.7 cm \\
                \textbf{Ours} & \textbf{25.9 cm} \\
                \bottomrule
            \end{tabular}
        \end{block}
        \begin{block}{UCF-Crime crime type (top-1)}
            \scriptsize
            Ours 43.10\%; Wu et al. CVPR 2024: 41.43\%; Anomize CVPR 2025: 47.14\% (both use CLIP).
        \end{block}
    \end{columns}
\end{frame}

% --- SLIDE 7: WORKING PROTOTYPE DEMONSTRATION ---
\begin{frame}{7. Working Prototype Demonstration}
    \framesubtitle{End-to-End Application Walkthrough \& Visual Evidence}
    \begin{columns}[T]
        \column{0.48\textwidth}
        \centering
        \figorbox{fig_scene1.png}{\linewidth}{Scene 1 map: figures/fig\_scene1.png}
        {\scriptsize Building, 30 min: layout and the map built by the robot}

        \column{0.48\textwidth}
        \centering
        \figorbox{fig_demo_output.png}{\linewidth}{Scene-understanding output: figures/fig\_demo\_output.png}
        {\scriptsize Scene understanding on an unseen video: crime type, kind of emergency, strongest moment}
    \end{columns}
    \vspace{0.15cm}

    \begin{block}{Live Demonstration Scenario}
        \footnotesize The robot starts in the unknown rescue field, follows the operator's route through the damaged building, maps the field and returns to its start point; a recorded video is then classified by the scene-understanding network.
    \end{block}
\end{frame}

% =========================================================================
% SECTION 4: TECHNICAL KNOWLEDGE & VIVA VOCE
% =========================================================================
\section{Technical Knowledge}

% --- SLIDE 8: TECHNICAL JUSTIFICATION & ARCHITECTURAL DECISIONS ---
\begin{frame}{8. Technical Justification \& Architectural Decisions}
    \framesubtitle{Platform Rationale, Engineering Trade-offs \& Operational Considerations}
    \begin{columns}[T]
        \column{0.48\textwidth}
        \begin{block}{1. Platform \& Technology Justification}
            \scriptsize
            \begin{itemize}
                \item One depth camera instead of lidar: cheaper, lighter, and the same Kinect gives colour video for scene understanding.
                \item CoppeliaSim gives exact ground truth for honest scoring; VideoMamba scales linearly with clip length, unlike attention.
            \end{itemize}
        \end{block}

        \begin{block}{2. Architecture \& Integration Trade-offs}
            \scriptsize
            \begin{itemize}
                \item A 2D slice of the depth image keeps SLAM fast but sees only 60 deg, so open floor relies on odometry.
                \item Agents share one process on a message bus: low latency, and each agent can be replaced alone.
            \end{itemize}
        \end{block}

        \column{0.48\textwidth}
        \begin{block}{3. Usability, Security \& Safety Measures}
            \scriptsize
            \begin{itemize}
                \item 0 collisions measured; 0.5 m standoff from unknown space; drop check at floor edges; stall watchdog.
                \item Ground truth used only for scoring; an unsure crime answer is shown as unsure.
            \end{itemize}
        \end{block}

        \begin{block}{4. Operational Constraints \& Future Scope}
            \scriptsize
            \begin{itemize}
                \item Navigation tested in simulation; error grows on long open floor; about one normal video in five is flagged.
                \item Next: wider view or visual odometry, more video for rare crimes, real-robot tests.
            \end{itemize}
        \end{block}
    \end{columns}
\end{frame}

% =========================================================================
% REFERENCES & THANK YOU
% =========================================================================
\section{References}
\nocite{*}
\begin{frame}[allowframebreaks]{References}
    \bibliographystyle{ieeetr}
    \bibliography{ref}
\end{frame}

\begin{frame}
    \centering
    {\Huge \textbf{\textcolor{amritaMaroon}{Thank You}}} \\[0.4cm]
    {\large \textbf{Questions \& Viva Voce}}
\end{frame}

\end{document}
